\section{Результаты}

\subsection{Описание разработанного программного обеспечения}

Разработано Spring Boot приложение, реализующее серверную часть для \textbf{шести доменов}, объединённых единым механизмом идемпотентности:
\begin{itemize}
    \item \textbf{Заказы} (\texttt{/api/orders}) --- аннотационный \texttt{@RestController};
    \item \textbf{Платежи} (\texttt{/api/payments}) --- функциональные эндпоинты (\texttt{RouterFunction});
    \item \textbf{Товары} (\texttt{/api/products}) --- автогенерируемые эндпоинты Spring Data REST;
    \item \textbf{Отправления} (\texttt{/api/soap/}) --- SOAP-эндпоинт на Spring-WS, контракт \texttt{xsd/shipments.xsd};
    \item \textbf{Уведомления} (\texttt{localhost:9090}) --- gRPC-сервис, контракт \texttt{notification.proto};
    \item \textbf{Инвентарь} (\texttt{/graphql}) --- GraphQL-контроллер, схема \texttt{schema.graphqls}.
\end{itemize}

\textbf{Стек технологий:} Spring Boot 3.3.6, Java 17, PostgreSQL 16, Flyway, Testcontainers, Docker Compose. Дополнительно: Spring-WS 4.0 + wsdl4j (SOAP), \texttt{grpc-spring-boot-starter 3.1} + grpc-java 1.65 + protobuf 3.25 (gRPC), \texttt{spring-boot-starter-graphql} (GraphQL).

\subsubsection{Структура проекта}

\begin{verbatim}
src/main/java/com/example/idempotency/
|-- config/                 # WebMvcConfig (URL-паттерны фильтра), IdempotencyProperties
|-- idempotency/            # Filter, Service, Entity, Repository, ResponseKind
|-- order/                  # @RestController + Service + Repository
|-- payment/                # Functional: PaymentRouter + Handler
|-- product/                # Spring Data REST: ProductRepository + EventHandler
|-- soap/shipment/          # SOAP: ShipmentEndpoint + Service + WebServiceConfig
|-- grpc/                   # GrpcIdempotencyInterceptor (общий для всех gRPC-сервисов)
|   `-- notification/       # gRPC-домен: NotificationGrpcService + Service + Entity
|-- graphql/inventory/      # GraphQL: InventoryGraphqlController + Service + Entity
+-- exception/              # Глобальный обработчик ошибок
src/main/proto/notification.proto
src/main/resources/
|-- xsd/shipments.xsd
|-- graphql/schema.graphqls
|-- db/migration/                    # V1..V7
+-- db/migration-non-idempotent/     # V3, V8
\end{verbatim}

\subsubsection{Схема базы данных}

База данных состоит из \textbf{семи таблиц}, управляемых миграциями Flyway V1--V7 (плюс V8 в профиле non-idempotent):

\begin{table}[H]
\centering
\caption{Схема базы данных}
\label{tab:db-schema}
\scriptsize
\begin{tabularx}{\textwidth}{@{}l X l l@{}}
\toprule
\textbf{Таблица} & \textbf{Ключевые поля} & \textbf{Constraints} & \textbf{Назначение} \\
\midrule
orders            & id (UUID), description, amount, status                                                    & UNIQUE(desc, amount)            & REST/@RestController \\
payments          & id (UUID), order\_id (FK), amount, status                                                  & FK $\to$ orders                 & REST/Functional \\
products          & id (BIGSERIAL), sku, name, price, version                                                  & UNIQUE(sku)                     & Spring Data REST \\
shipments         & id (UUID), tracking\_number, recipient, status                                             & UNIQUE(tracking\_number)        & SOAP \\
notifications     & id (UUID), recipient, message, message\_hash, status                                       & UNIQUE(recipient, msg\_hash)    & gRPC \\
inventories       & id (UUID), item\_name, quantity                                                            & UNIQUE(item\_name)              & GraphQL \\
idempotency\_keys & id, key, method, path, response\_status, response\_body, response\_content\_type, response\_kind, response\_body\_bytes & UNIQUE(key) & Снэпшоты ответов (V4) \\
\bottomrule
\end{tabularx}
\end{table}

Расширение \texttt{idempotency\_keys} в миграции V4 (\S\,\ref{subsec:transport-agnostic-snapshot}) добавляет три колонки, без которых фильтр не мог бы корректно реплеить SOAP/XML и gRPC/protobuf. Поле \texttt{response\_kind} различает текстовое тело (\texttt{TEXT}, используется JSON/XML) и бинарное (\texttt{BINARY}, используется protobuf); \texttt{response\_content\_type} восстанавливает оригинальный MIME-тип на реплее. TTL ключей задаётся колонкой \texttt{expires\_at} (по умолчанию 24 часа).

\subsubsection{Компоненты механизма идемпотентности}

Механизм идемпотентности реализован через четыре компонента, переиспользуемых четырьмя HTTP-транспортами и одним gRPC-перехватчиком:

\begin{enumerate}
    \item \textbf{IdempotencyFilter} --- единственный HTTP-фильтр (наследник \mbox{\texttt{OncePerRequestFilter}}), зарегистрированный на URL-паттернах четырёх транспортов: \texttt{/api/orders/*}, \texttt{/api/payments/*}, \texttt{/api/soap/*}, \texttt{/graphql}. Реализует алгоритм~\ref{alg:idempotency-check}. Отключается параметром \mbox{\texttt{app.idempotency.enabled=false}}. После доработки V4 (\S\,\ref{subsec:transport-agnostic-snapshot}) транспорт-агностичен.

    \item \textbf{GrpcIdempotencyInterceptor} --- \texttt{@GrpcGlobalServerInterceptor}, делегирующий тому же \texttt{IdempotencyService}. Generic-реализация работает с любым unary-методом любого gRPC-сервиса через \texttt{MethodDescriptor.getResponseMarshaller()}.

    \item \textbf{IdempotencyService} --- единая точка интеграции для всех транспортов. Методы \texttt{findByKey()}, \texttt{registerKey()}, \texttt{saveTextResponse()} (для JSON/XML), \texttt{saveBinaryResponse()} (для protobuf-байтов). Операции транзакционны.

    \item \textbf{IdempotencyKeyEntity} --- JPA-сущность, представляющая запись в таблице \texttt{idempotency\_keys}. Содержит ключ, HTTP-метод, путь запроса, снэпшот ответа ($\langle$status, contentType, kind, body$\rangle$) и временные метки.
\end{enumerate}

\subsubsection{Профили конфигурации}

Приложение поддерживает два режима работы через Spring Profiles:

\begin{itemize}
    \item \textbf{Профиль по умолчанию}: идемпотентность включена (\mbox{\texttt{enabled=true}}), Flyway применяет миграции V1 и V2, все constraints активны.

    \item \textbf{Профиль non-idempotent}: идемпотентность отключена (\mbox{\texttt{enabled=false}}), Flyway дополнительно применяет миграции V3 (дроп UNIQUE на \texttt{orders}, \texttt{products}) и V8 (дроп UNIQUE на \texttt{shipments}, \texttt{notifications}, \texttt{inventories}). gRPC-перехватчик так же проверяет флаг \texttt{enabled} и при выключенном режиме пропускает вызов без перехвата.
\end{itemize}

Переключение профилей позволяет запускать одни и те же тесты в двух конфигурациях и наглядно демонстрировать последствия отсутствия механизмов идемпотентности на всех шести транспортах.

\subsection{Результаты тестирования (профиль idempotent)}

Полный набор из 36 интеграционных тестов проходит успешно:

\begin{table}[H]
\centering
\caption{Результаты тестирования (объединённый прогон обоих профилей)}
\label{tab:results-idempotent}
\small
\begin{tabular}{@{}lccl@{}}
\toprule
\textbf{Тестовый класс} & \textbf{Тестов} & \textbf{Результат} & \textbf{Тип контроллера / транспорт} \\
\midrule
OrderIdempotencyTest       & 6  & PASS & @RestController (HTTP/JSON) \\
PaymentIdempotencyTest     & 5  & PASS & Functional Endpoints (HTTP/JSON) \\
ProductIdempotencyTest     & 5  & PASS & Spring Data REST (HTTP/HAL) \\
SoapIdempotencyTest        & 3  & PASS & Spring-WS (SOAP/XML) \\
GrpcIdempotencyTest        & 4  & PASS & @GrpcService (HTTP/2 + protobuf) \\
GraphqlIdempotencyTest     & 4  & PASS & spring-boot-starter-graphql \\
DataImportTest             & 2  & PASS & Импорт JSON/CSV \\
NonIdempotentProfileTest   & 7  & PASS & Демонстрация багов на 6 транспортах \\
\midrule
\textbf{Итого}             & \textbf{36} & \textbf{PASS} & \\
\bottomrule
\end{tabular}
\end{table}

\subsection{Тест-кейсы по контроллерам}

\subsubsection{@RestController (OrderIdempotencyTest)}

\begin{enumerate}
    \item \textbf{POST с одинаковым Idempotency-Key дважды} --- создаётся только один заказ, второй запрос возвращает кэшированный ответ. Проверяется через утверждение \mbox{\texttt{assertEquals(1, count)}}.
    \item \textbf{POST с разными ключами} --- создаются два разных заказа, что подтверждает: идемпотентность привязана к конкретному ключу, а не к содержимому запроса.
    \item \textbf{POST без Idempotency-Key} --- возвращается 400 Bad Request, что предотвращает случайное создание неидемпотентных запросов.
    \item \textbf{PUT идемпотентен} --- повторное обновление с теми же данными не изменяет состояние. Проверяется сравнением полей \texttt{description}, \texttt{amount}, \texttt{status} после двух последовательных PUT.
    \item \textbf{DELETE идемпотентен} --- первый вызов возвращает 204 No Content (ресурс удалён), второй --- 404 Not Found (ресурс уже отсутствует). Оба варианта соответствуют идемпотентному поведению.
    \item \textbf{GET идемпотентен} --- три последовательных GET-запроса возвращают одинаковый JSON-ответ.
\end{enumerate}

\subsubsection{Functional Endpoints (PaymentIdempotencyTest)}

\begin{enumerate}
    \item \textbf{POST платежа с одним ключом дважды} --- создаётся один платёж. Демонстрирует, что IdempotencyFilter работает с RouterFunction так же, как с @RestController.
    \item \textbf{POST с разными ключами} --- создаются два платежа.
    \item \textbf{POST без ключа} --- 400 Bad Request.
    \item \textbf{GET идемпотентен} --- стабильный результат при многократных запросах.
    \item \textbf{Retry (3 запроса с одним ключом)} --- создаётся один платёж, несмотря на 3 попытки. Это ключевой production-сценарий: клиент повторяет запрос из-за таймаута, но система обрабатывает его только один раз.
\end{enumerate}

\subsubsection{Spring Data REST (ProductIdempotencyTest)}

\begin{enumerate}
    \item \textbf{POST дубликата SKU} --- 409 Conflict. UNIQUE constraint на поле \texttt{sku} предотвращает создание товара с тем же артикулом.
    \item \textbf{PUT с корректной версией} --- успешное обновление, \texttt{@Version} автоматически инкрементируется. Проверяется совпадение \mbox{\texttt{version~+~1}} с актуальной версией.
    \item \textbf{Concurrent PUT (два обновления со старой версией)} --- первое обновление проходит успешно (версия 0 $\to$ 1), второе получает 412 Precondition Failed, поскольку передан устаревший ETag. Optimistic locking предотвращает lost update.
    \item \textbf{PATCH идемпотентен} --- повторное применение одного и того же PATCH не изменяет данные (имя товара остаётся тем же после второго PATCH).
    \item \textbf{GET идемпотентен} --- стабильный результат.
\end{enumerate}

\subsubsection{Импорт данных (DataImportTest)}

\begin{enumerate}
    \item \textbf{Импорт заказов из JSON} --- загрузка данных из файла \mbox{\texttt{orders-input.json}}, создание через API, сравнение количества, статусов и сумм с эталонным файлом.
    \item \textbf{Импорт товаров из CSV} --- загрузка данных из файла \mbox{\texttt{products-input.csv}} через Jackson CSV, сохранение в БД, затем повторная попытка импорта: все 3 попытки отклоняются UNIQUE constraint, количество записей не увеличивается.
\end{enumerate}

\subsubsection{SOAP (SoapIdempotencyTest)}

Тесты используют \texttt{TestRestTemplate} на random-порту (MockMvc маршрутизирует только в \texttt{DispatcherServlet} и не вызывает \texttt{MessageDispatcherServlet}), отправляя сырой SOAP-envelope в \texttt{/api/soap/}:
\begin{enumerate}
    \item \textbf{Один Idempotency-Key дважды} --- ответы байт-в-байт идентичны (включая Content-Type \texttt{text/xml}), в \texttt{shipments} ровно одна запись. Это подтверждает, что доработанный фильтр (\S\,\ref{subsec:transport-agnostic-snapshot}) корректно реплеит XML-тело.
    \item \textbf{Разные ключи} --- две записи \texttt{shipments}.
    \item \textbf{POST без Idempotency-Key} --- 400 Bad Request, ничего не создано.
\end{enumerate}

\subsubsection{gRPC (GrpcIdempotencyTest)}

Тесты подключаются через \texttt{InProcessChannelBuilder} к in-process gRPC-серверу (\texttt{grpc.server.in-process-name=idempotency-test}); идемпотентный ключ передаётся в \texttt{Metadata}:
\begin{enumerate}
    \item \textbf{\texttt{Send} с одним ключом дважды} --- \texttt{SendResponse.id} совпадает, в \texttt{notifications} одна запись. Реплей происходит через десериализацию байтов из \texttt{response\_body\_bytes} тем же marshaller, что использовался при сериализации.
    \item \textbf{Разные ключи} --- две записи \texttt{notifications}.
    \item \textbf{Без \texttt{idempotency-key}} --- \texttt{Status.INVALID\_ARGUMENT}, ничего не создано.
    \item \textbf{Бинарная идентичность реплея} --- \texttt{toByteArray()} обоих ответов совпадают побайтно. Это сильнее, чем равенство объектов: подтверждает, что сериализованная форма стабильна.
\end{enumerate}

\subsubsection{GraphQL (GraphqlIdempotencyTest)}

Тесты используют \texttt{TestRestTemplate} на \texttt{/graphql}, отправляя GraphQL-запросы JSON-обёрткой:
\begin{enumerate}
    \item \textbf{Мутация с одним Idempotency-Key дважды} --- JSON-ответы идентичны, в \texttt{inventories} одна запись.
    \item \textbf{Мутации с разными ключами} --- две записи.
    \item \textbf{Мутация без ключа} --- 400 Bad Request.
    \item \textbf{Query после мутации} --- запрос возвращает только что созданную запись, что демонстрирует совместимость GraphQL Query с механизмом идемпотентности (Query через POST требует ключ, но это не мешает выборке данных).
\end{enumerate}

\subsection{Результаты профиля non-idempotent}

Профиль \texttt{non-idempotent} отключает механизмы идемпотентности на всех слоях:
\begin{itemize}
    \item \texttt{IdempotencyFilter} пропускает все POST-запросы без проверки заголовка (\texttt{enabled=false});
    \item \texttt{GrpcIdempotencyInterceptor} обходит логику регистрации ключа по тому же флагу;
    \item Миграция V3 удаляет UNIQUE constraints на \texttt{orders} и \texttt{products};
    \item Миграция V8 удаляет UNIQUE constraints на \texttt{shipments}, \texttt{notifications}, \texttt{inventories}.
\end{itemize}

\begin{table}[H]
\centering
\caption{Демонстрация проблем без идемпотентности (профиль non-idempotent, все 6 транспортов)}
\label{tab:results-non-idempotent}
\small
\begin{tabularx}{\textwidth}{@{}X X c@{}}
\toprule
\textbf{Тест} & \textbf{Обнаруженная проблема} & \textbf{Статус} \\
\midrule
REST POST без Idempotency-Key       & Запрос проходит без проверки --- любой POST создаёт ресурс         & BUG \\
REST повторный POST с тем же ключом & Создаётся дубликат заказа --- ключ игнорируется                    & BUG \\
Data REST дубликат SKU              & Создаётся второй товар с тем же артикулом                          & BUG \\
Functional retry платежа (3 раза)   & 3 платежа вместо 1 --- тройное списание средств                    & BUG \\
SOAP createShipment с тем же tracking\_number & Создаются два Shipment c одинаковым номером                & BUG \\
gRPC Send с тем же idempotency-key  & Interceptor выключен $\Rightarrow$ дубликат уведомления            & BUG \\
GraphQL мутация с тем же ключом     & Фильтр выключен $\Rightarrow$ два Inventory с одним item\_name     & BUG \\
\bottomrule
\end{tabularx}
\end{table}

Все семь тестов \texttt{NonIdempotentProfileTest} проходят зелёными, подтверждая наличие проблем во всех шести транспортах. Каждый тест использует прямые assertions: например, проверяется, что \mbox{\texttt{notificationRepository.count() == 2}} --- без перехватчика gRPC и без DB constraint действительно создаются дубликаты. Это демонстрирует критическое значение комбинации \emph{транспортного механизма} (фильтр/перехватчик) и \emph{domain-уровня} (UNIQUE constraint) для корректного функционирования системы независимо от выбранного протокола.
